% gerado por ferramentas/md2tex_enunciado.py a partir de questoes/q121.md (não editar à mão)
Um estudante acende uma lâmpada de 6\,W de luz vermelha (comprimento de onda igual a 650\,nm) em um quarto escuro. A 1,0\,m de distância dessa lâmpada, encontra-se uma planta cujas folhas têm área de $4{,}0\,\text{cm}^2$. Em uma rápida pesquisa na internet, o estudante verifica que plantas dessa mesma espécie possuem espectro de absorção conforme ilustrado no gráfico, que apresenta a porcentagem de luz incidente absorvida em função do comprimento de onda.\par
\figq[0.44\linewidth]{q121_fig1.png}
Considere a lâmpada como uma fonte puntiforme cuja intensidade luminosa é dada pela razão entre a sua potência e a área da superfície esférica formada a uma distância \textit{r} da fonte.\par
Utilize 3 como valor aproximado para $\pi$.\par
Com base nessas informações e no gráfico, qual é o valor mais próximo da energia, em joule, que cada folha absorverá a cada segundo?\par
\begin{altlinha}
\alth{A}{$3{,}8 \times 10^{-4}$\,J}%
\alth{B}{$4{,}8 \times 10^{-4}$\,J}%
\alth{C}{$2{,}0 \times 10^{-4}$\,J}%
\alth{D}{$1{,}8 \times 10^{-4}$\,J}%
\alth{E}{$1{,}2 \times 10^{-4}$\,J}
\end{altlinha}
